\capitulo{8}{Lineas de trabajo futuras}


Como continuación del presente trabajo, se propone el desarrollo de una línea de investigación orientada a la evolución del sistema hacia una plataforma híbrida que combine análisis cualitativo y cuantitativo del movimiento. En este sentido, uno de los principales avances consistiría en la incorporación de un módulo de validación cuantitativa basado en la obtención de parámetros numéricos objetivos derivados de los datos capturados por el guante sensorizado.

Dicho módulo permitiría transformar las señales recogidas (posición, velocidad, aceleración y orientación) en indicadores medibles, tales como amplitud de movimiento, velocidad angular, frecuencia de ejecución, estabilidad y simetría entre extremidades. Estos parámetros podrían calcularse mediante la aplicación de modelos matemáticos y físicos básicos, como derivadas temporales para el cálculo de velocidad y aceleración, o métricas estadísticas para evaluar la variabilidad del movimiento. La integración de estos valores facilitaría la comparación con rangos de normalidad previamente definidos en función de la edad y el desarrollo motor del paciente, contribuyendo así a una evaluación más objetiva y reproducible.

Además, la incorporación de técnicas de análisis de datos y aprendizaje automático permitiría identificar patrones característicos asociados a posibles alteraciones neuromotoras, mejorando la capacidad del sistema para apoyar la detección temprana. En esta línea, resulta especialmente relevante el desarrollo de modelos personalizados que se adapten a las características individuales de cada paciente, teniendo en cuenta su evolución a lo largo del tiempo. El uso de datos longitudinales permitiría no solo detectar desviaciones puntuales, sino también analizar tendencias y progresión de la patología, aportando un valor significativo en el seguimiento clínico.

Por otra parte, se plantea como mejora del dispositivo hardware el diseño de un guante inteligente que integre una pantalla o interfaz visual embebida. Esta pantalla permitiría proporcionar retroalimentación en tiempo real tanto a profesionales sanitarios como a los propios usuarios. A nivel clínico, se podrían mostrar métricas relevantes, alertas ante desviaciones significativas y representaciones gráficas simplificadas de la evolución del paciente. Desde la perspectiva del usuario, especialmente en el caso de niños, la interfaz podría ofrecer indicaciones visuales intuitivas, refuerzo positivo y elementos de gamificación que favorezcan la adherencia al uso del sistema.

Adicionalmente, se propone la integración del sistema con plataformas de telemedicina, lo que permitiría la monitorización remota de los pacientes y el acceso a los datos por parte de profesionales sanitarios en tiempo real. Esta funcionalidad facilitaría la continuidad asistencial, reduciría la necesidad de desplazamientos y ampliaría el alcance del sistema a entornos domiciliarios o rurales.

Finalmente, la combinación de estos avances abriría la puerta a nuevas aplicaciones, como el seguimiento longitudinal de pacientes, la monitorización remota y el uso del sistema como herramienta de apoyo en procesos de rehabilitación. En conjunto, estas líneas de investigación contribuirían a mejorar la precisión, usabilidad y aplicabilidad clínica del sistema propuesto, consolidándolo como una herramienta innovadora en el ámbito de la evaluación y apoyo al diagnóstico de alteraciones neuromotoras.